\documentclass[11pt]{article}
\title{On the Closure of Compositional Hierarchies in Digital Symmetries}
\author{A rigorous researcher-engineer bridging abstract mathematics and concrete digital systems.}
\usepackage[margin=1in]{geometry}
\usepackage{graphicx}
\usepackage{amsmath,amssymb,mathtools}
\usepackage{hyperref}
\graphicspath{{media/}{media/diagrams/}{images/}{artwork/}}

\begin{document}

\section*{Cryptographic Constructions \& Symmetry Constraints}

This chapter examines how symmetry constraints shape both the design and the analysis of cryptographic primitives. The central theme is that useful cryptographic behavior often requires a careful balance: enough structure to support efficient implementation and systematic reasoning, but enough asymmetry to resist unintended invariants, equivalences, and attacks. In this setting, symmetry is not merely an aesthetic property; it is a design parameter that can strengthen or weaken security depending on how it interacts with composition.

A recurring quantitative lens is the round-dependent error profile \(\varepsilon(r)\), where \(r\) denotes the number of rounds in a construction. In practice, \(\varepsilon(r)\) may summarize how quickly avalanche behavior emerges, or how differential and linear biases decay as rounds are added. From the standpoint of closure and compositional hierarchies, the key question is whether the intended mixing behavior is preserved under iteration, or whether hidden invariants survive round composition and create exploitable regularities.

For substitution--permutation networks and related designs, the S-box is the primary local source of nonlinearity. Its symmetry profile is therefore critical. Two S-boxes related by affine equivalence or CCZ equivalence may share important cryptanalytic properties, even when their concrete representations differ. In particular, nonlinearity and differential uniformity are standard measures used to assess resistance to linear and differential attacks, and they also serve as coarse invariants under many symmetry-preserving transformations. The design problem is not simply to maximize these metrics in isolation, but to understand which equivalence classes preserve them and which transformations alter the effective attack surface.

Round composition introduces a second layer of structure. In SPN and Feistel constructions, the interaction between local nonlinear layers, linear diffusion, and key mixing can create or destroy invariants across rounds. A single round may exhibit obvious symmetries, but repeated composition can either break them or propagate them in disguised form. This makes round-reduced experiments especially useful: by studying truncated versions of a cipher, one can often detect persistent algebraic relations, invariant subspaces, or symmetry-induced biases that are difficult to observe directly in the full construction. Such experiments do not prove insecurity by themselves, but they provide evidence about how closure behaves under iteration.

The chapter also treats group actions as a unifying language for design and cryptanalysis. A group action may encode input/output relabelings, affine transformations, coordinate permutations, or other admissible symmetries of a primitive. From the design perspective, these actions help classify constructions up to natural equivalence. From the cryptanalytic perspective, they can reveal when two apparently different instances are actually the same object in disguise, or when a family of candidates collapses into a smaller number of symmetry classes. The main pitfall is closure: if the set of allowed transformations is not carefully controlled, then an analysis may inadvertently identify non-equivalent objects, or a design may inherit unintended invariants from a larger symmetry group than intended.

Accordingly, the chapter emphasizes a practical distinction between symmetry as a tool and symmetry as a vulnerability. In some contexts, symmetry supports efficient implementation, compact specification, or systematic search over design spaces. In others, it creates fixed points, invariant partitions, or low-dimensional orbits that can be exploited by an attacker. The relevant question is not whether a construction has symmetries, but which symmetries are admissible, which are broken by the round function, and which survive composition in a way that undermines diffusion or confusion.

The accompanying artifact is an equivariance classifier intended to support this analysis. Its role is to categorize candidate transformations and identify whether a given primitive, component, or reduced-round instance is invariant, equivariant, or neither under a specified action. In parallel, round-reduced experiments provide empirical evidence about the persistence of symmetry-related phenomena across a small number of rounds. These artifacts are meant to be reproducible and to support comparison across design candidates rather than to replace formal proofs.

Because cryptographic evaluation can have real-world consequences, this chapter also includes an ethics note. Symmetry analysis and round-reduction experiments should be used responsibly, with attention to disclosure norms, reproducibility, and the limits of empirical evidence. The goal is to improve understanding of cryptographic structure and to support safer design, not to encourage misuse of analytic techniques against deployed systems.

Overall, the chapter situates cryptographic constructions within the book's broader framework of closure, composition, and symmetry. It shows that the same conceptual tools used to study digital hierarchies in general also illuminate the fine structure of S-boxes, round functions, and equivalence relations in modern cryptography.

\paragraph{Section artifact.} The section's executable companion is an equivariance classifier for candidate transformations together with a small suite of round-reduced experiments. The intended workflow is to test a primitive against a specified action, record whether the result is invariant, equivariant, or neither, and then compare those outcomes across reduced-round variants to estimate how quickly symmetry-induced structure decays.

\paragraph{Interpretive note.} The chapter uses symmetry language in a deliberately conservative way. Equivalence notions such as affine and CCZ equivalence are treated as classification tools, not as security guarantees. Likewise, empirical regularities observed in reduced-round settings are treated as indicators for further analysis, not as standalone proofs of attack feasibility.

\paragraph{Ethics note.} Any implementation of the artifact should be used in accordance with responsible disclosure practices and local policy. The purpose of the analysis is to support robust cryptographic design, reproducible evaluation, and careful interpretation of symmetry effects, especially when empirical findings may be sensitive for deployed systems.

\paragraph{Technical focus.} The main technical objects in this chapter are the round function, the S-box equivalence class, and the induced action of admissible symmetries on the space of candidate constructions. The chapter uses these objects to organize three related questions: how quickly \(\varepsilon(r)\) decays with round count, which equivalence relations preserve the relevant cryptanalytic metrics, and which round-composition laws preserve or destroy invariants.

\paragraph{Design and analysis workflow.} A typical workflow is to begin with a candidate S-box or round function, compute its nonlinearity and differential uniformity, and then compare these values across affine or CCZ-equivalent representatives. One then embeds the component into an SPN or Feistel template, studies the effect of round composition on avalanche and bias, and checks whether any invariant partitions or low-dimensional orbits remain visible under the chosen group action. The resulting evidence can guide both design refinement and cryptanalytic prioritization.

\paragraph{Caution on closure.} Symmetry-preserving transformations can be useful for classification, but they can also obscure distinctions that matter for security. In particular, closure under a large transformation family may collapse a design space into a small number of orbits, making it easier to search but also easier for an attacker to exploit shared structure. The chapter therefore treats closure as a property to be measured, not assumed.

\paragraph{Summary.} In short, cryptographic constructions live at the intersection of nonlinearity, composition, and symmetry. The chapter's contribution is to make that intersection explicit: to show how S-box symmetries, round composition, and group actions interact, and to provide an artifact-driven method for tracking when those interactions are benign, informative, or dangerous.


\end{document}
